
\section{Background}

\subsection{Softmax Attention}

Standard scaled dot-product attention:

\[
\text{Attn}(Q, K, V) = \text{softmax}\!\left(\frac{QK^\top}{\sqrt{d}}\right)V
\]

\textbf{Properties:}
\begin{itemize}
\item \textbf{Non-negativity:} $\text{softmax}(\cdot) > 0$ elementwise.
\item \textbf{Normalization:} $\sum_j \text{softmax}(\cdot)_{ij} = 1$.
\item \textbf{Locality:} With scale factor $1/\sqrt{d}$, softmax concentrates on
  the largest entries of $QK^\top$.

\end{itemize}
\subsection{Linear Attention}

Replacing softmax with a kernel factorization:

\[
\exp(q^\top k) \approx \phi(q)^\top \phi(k), \quad \phi: \mathbb{R}^d \to \mathbb{R}^m
\]

allows rewriting attention as:

\[
\text{Attn}(Q_i, K, V) = \frac{\sum_j \phi(q_i)^\top \phi(k_j) V_j}{\sum_j \phi(q_i)^\top \phi(k_j)} = \frac{\phi(q_i)^\top \sum_j \phi(k_j) V_j}{\phi(q_i)^\top \sum_j \phi(k_j)}
\]

\textbf{Complexity.} The sums $\sum_j \phi(k_j) V_j$ and $\sum_j \phi(k_j)$ are 
$\mathcal{O}(L m d)$, independent of the query position. Precomputed once, 
each query costs $\mathcal{O}(m d)$.

\textbf{Common kernels:}
\begin{itemize}
\item \textbf{ELU + 1} (Katharopoulos et al., 2020): $\phi(x) = \text{elu}(x) + 1$.
\item \textbf{ReLU:} $\phi(x) = \text{ReLU}(x)$.
\item \textbf{Softplus:} $\phi(x) = \log(1 + e^x)$.
\item \textbf{Random Fourier Features:} $\phi(x) = e^{\omega^\top x}$ for random
  $\omega$.

\end{itemize}
All are non-negative; the first three are cheap and deterministic.

\subsection{Locality and Its Loss}

Softmax attention exhibits locality: nearby tokens receive higher attention 
mass on average, especially with positional information. This is critical 
for language understanding.

\textbf{Why linear attention loses locality.} The kernel approximation is 
position-independent: $\phi(q_i)$ and $\phi(k_j)$ have no information about 
$|i - j|$. So there is no bias toward nearby tokens.

\textbf{Empirical effect.} On tasks requiring local context (NER, chunked 
summarization), linear attention underperforms softmax.

\subsection{Cosformer's Reweighting}

Cosformer adds cosine/sine modulation:

\[
\tilde{\phi}(q_i) = \phi(q_i) \cdot \cos\!\left(\frac{\pi i}{2 L}\right), \quad
\tilde{\phi}'(q_i) = \phi(q_i) \cdot \sin\!\left(\frac{\pi i}{2 L}\right)
\]

Then the attention becomes:

\[
\text{Attn}(Q_i, K, V) = \frac{\tilde{\phi}(q_i)^\top \sum_j \tilde{\phi}(k_j) V_j + \tilde{\phi}'(q_i)^\top \sum_j \tilde{\phi}'(k_j) V_j}{\tilde{\phi}(q_i)^\top \sum_j \tilde{\phi}(k_j) + \tilde{\phi}'(q_i)^\top \sum_j \tilde{\phi}'(k_j)}
\]

\textbf{Why it works.} The identity:

\[
\cos(\theta_i - \theta_j) = \cos\theta_i \cos\theta_j + \sin\theta_i \sin\theta_j
\]

means the sum $\tilde{\phi}(q_i)^\top \tilde{\phi}(k_j) + \tilde{\phi}'(q_i)^\top \tilde{\phi}'(k_j)$ 
equals $\phi(q_i)^\top \phi(k_j) \cdot \cos(\theta_i - \theta_j)$. The 
cosine decays with $|i - j|$, giving locality.

\subsection{GQA and RoPE}

\textbf{GQA.} $H_{kv} < H$ reduces KV cache. Handled via \texttt{repeat\_kv}.

\textbf{RoPE.} Rotation of $Q, K$ by position-dependent angles. Orthogonal to 
Cosformer's reweighting; both can coexist.

\subsection{Notation Summary}

\begin{table}[H]
\centering
\small
\begin{tabular}{ll}
\toprule
\textbf{Symbol} & \textbf{Meaning} \\
\midrule
$L$ & Sequence length \\
$d$ & Head dimension \\
$H$ & Query head count \\
$H_{kv}$ & KV head count \\
$\phi$ & Kernel feature map \\
$\theta_i$ & Position angle ($\pi i / 2L$) \\
$r$ & LoRA rank \\
$C$ & Chunk size \\
\bottomrule
\end{tabular}
\end{table}
</parameter>
\begin{lstlisting}

---
\end{lstlisting}